% ChargeShield-FL --- DSN 2027 --- Figura: i tre punti di osservazione
% dell'avversario (2026-09-16, Sprint 10zz+117).
%
% NOTA TIKZ: ogni \\ deve restare al livello ESTERNO del testo del nodo. Un
% \\ annidato dentro l'argomento di un altro comando (es. \emph{a\\b}) rompe
% il layout di align con "Undefined control sequence \pgfutil@next" --- bug
% reale trovato e corretto nel Sprint 10zz+109 su ml_plane_architecture.tex.
\begin{figure}[t]
\centering
\begin{tikzpicture}[
 font=\footnotesize,
  stage/.style={draw, rounded corners, minimum height=0.65cm,
                minimum width=1.55cm, align=center, font=\scriptsize},
  zone/.style={draw, dashed, rounded corners, inner sep=0.15cm},
  adv/.style={draw, circle, minimum size=0.48cm, thick,
              red!65!black, font=\scriptsize\bfseries},
  arr/.style={-Latex, thick}
]

  \node[stage] (raw)   {local update\\$g$};
  \node[stage, right=0.55cm of raw]  (clip)  {clipped\\$\bar{g}$};
  \node[stage, right=0.55cm of clip] (noise) {noised\\$\tilde{g}$};
  \node[stage, right=0.55cm of noise, minimum width=1.4cm] (agg) {aggregator\\FedAvg};

  \draw[arr] (raw)   -- node[above, font=\scriptsize]{clip} (clip);
  \draw[arr] (clip)  -- node[above, font=\scriptsize]{$+\,\mathcal{N}$} (noise);
  \draw[arr] (noise) -- node[above, font=\scriptsize]{send} (agg);

  % zona che non lascia mai il client
  \begin{scope}[on background layer]
    \node[zone, fit=(raw)(clip), label={[font=\scriptsize\itshape]below:client side}] {};
  \end{scope}

  \node[adv, above=0.35cm of raw]   (a1) {A1};
  \node[adv, above=0.35cm of clip]  (a2) {A2};
  \node[adv, above=0.35cm of noise] (a3) {A3};
  \draw[arr, red!65!black] (a1) -- (raw);
  \draw[arr, red!65!black] (a2) -- (clip);
  \draw[arr, red!65!black] (a3) -- (noise);

  \node[font=\scriptsize, align=center, below=0.55cm of raw]   {\texttt{dp-fedavg}\\upper bound\\\emph{beyond} S1};
  \node[font=\scriptsize, align=center, below=0.55cm of clip]  {\texttt{central}\\S1 aggregator};
  \node[font=\scriptsize, align=center, below=0.55cm of noise] {\texttt{local}\\S1 aggregator\\or eavesdropper};

\end{tikzpicture}
\caption{The three observation points are three adversary privileges on one
privatisation pipeline, not three privacy mechanisms. \textbf{A3} sees only
the clipped-and-noised update: an honest-but-curious aggregator, or a network
observer, under a local-DP deployment. \textbf{A2} sees the update after
clipping but before noise, which is what an aggregator legitimately receives
when the noise is added centrally --- still Scenario~1. \textbf{A1} sees the
raw update, which under our deployments never leaves the client
unprivatised; it therefore models an adversary \emph{strictly stronger} than
Scenario~1, and we evaluate it precisely because a null there bounds the
other two (\S\ref{sec:threat-observation}).}
\label{fig:adversary-points}
\end{figure}
